% =====================================================================
%  Namespaces Without a Name Server  --- draft 2 (pass one)
%  L. G. Meredith, F1R3FLY.io
% =====================================================================
\documentclass[11pt]{article}

\usepackage[T1]{fontenc}
\usepackage[utf8]{inputenc}
\usepackage{amsmath,amssymb,amsthm}
\usepackage[dvipsnames]{xcolor}
\usepackage[margin=1.1in]{geometry}
\usepackage{mathrsfs}
\usepackage{enumitem}
\usepackage{booktabs}
\usepackage{listings}
\usepackage[expansion=false]{microtype}
\usepackage[hidelinks]{hyperref}

\theoremstyle{definition}
\newtheorem{definition}{Definition}[section]
\newtheorem{construction}[definition]{Construction}
\newtheorem{notation}[definition]{Notation}
\newtheorem{convention}[definition]{Convention}
\theoremstyle{plain}
\newtheorem{proposition}[definition]{Proposition}
\newtheorem{lemma}[definition]{Lemma}
\newtheorem{theorem}[definition]{Theorem}
\newtheorem{corollary}[definition]{Corollary}
\theoremstyle{remark}
\newtheorem{remark}[definition]{Remark}
\newtheorem{example}[definition]{Example}

\newcommand{\rhoc}{\mathrm{rho}}
\newcommand{\quo}[1]{@#1}
\newcommand{\drp}[1]{{*}#1}
\newcommand{\inp}[3]{\mathsf{for}(\,#1 \leftarrow #2\,)#3}
\newcommand{\inpp}[5]{\mathsf{for}(\,#1 \leftarrow #2\ \&\ #3 \leftarrow #4\,)#5}
\newcommand{\out}[2]{#1!(#2)}
\newcommand{\enc}[1]{[\![#1]\!]}
\newcommand{\nequiv}{\equiv_{\mathcal{N}}}
\newcommand{\nil}{\mathbf{0}}
\newcommand{\Lft}{\mathsf{L}}
\newcommand{\Rgt}{\mathsf{R}}
\newcommand{\Frsh}{\mathsf{Fresh}_2}
\newcommand{\Dup}{\mathsf{D}}
\newcommand{\Ctx}{\mathscr{C}}
\newcommand{\Obs}{\mathcal{D}}
\newcommand{\bisim}{\approx}
\newcommand{\QD}{\mathrm{qd}}
\newcommand{\real}{\Vdash}
\newcommand{\dom}{\mathrm{dom}\,}
\newcommand{\fn}{\mathrm{fn}}
\newcommand{\nms}{\mathcal{N}\!\mathit{ms}}
\providecommand{\xLongrightarrow}[1]{\mathrel{\overset{#1}{\Longrightarrow}}}

\lstset{basicstyle=\ttfamily\small,columns=fullflexible,keepspaces=true,
        breaklines=true,frame=single,framesep=4pt,xleftmargin=6pt}

\title{\textbf{Namespaces Without a Name Server}\\[4pt]
\large Freshness by communication, and a fully abstract encoding of the
$\pi$-calculus in the rho calculus}
\author{L.\,G.\ Meredith\thanks{F1R3FLY.io.\ \texttt{lgreg.meredith@f1r3fly.io}}}
\date{August 2026}

\begin{document}
\maketitle

\begin{abstract}
\noindent
The rho calculus has no name-restriction operator and no ambient supply
of atoms: names are quoted processes. An encoding of the $\pi$-calculus
into it is therefore a test of a specific design objective --- that $\nu$
can be discharged without positing a runtime source of fresh names. The
encoding published with the calculus in 2005 was intended to meet that
objective and does not: Lybech exhibited two counterexamples and gave a
corrected encoding in which freshness is supplied by a central name
server reached on a fixed channel. That encoding is correct. It meets a
different design objective, and Lybech says as much in his conclusion,
where he describes a per-replication local server as closer to the
original intention and leaves it open on account of a scaffolding
problem: to build a source of fresh names one must already have one.

We settle the scaffolding problem. Three facts about the rho calculus do
the work. First, a reduction adds at most one name to a term, so
producing $n$ distinct names costs $n$ rendezvous, and no translation can
write them. Second, substitution reaches the object of a lift and does
not reach beneath a quote, so a name derived from a runtime parameter
must be built by communication --- one quote level per rendezvous.
Third, the rho duplicator reproduces quoted code verbatim, so copies of a
replicated body cannot be told apart by their syntax. Together these
force address arithmetic out of the translation and into a process, and
they determine that process almost uniquely: a two-communication
allocator we call $\Frsh$.

The architecture $\Frsh$ realises is Abramsky's. His proof expressions for
classical linear logic carry an invariant --- names are a root together
with a binary address, and addresses sharing a root are incompatible in
the prefix ordering --- which is exactly the invariant the 2005 encoding
required and never stated, and which its splitting of parameters across
parallel composition already maintains. What it omitted is the other half
of the Copy rule: the clause that re-contracts a copied box's side ports
so that the parent's namespace is never re-rooted. Abramsky's Copy cannot
be transcribed, because it renames the contents of a box at the metalevel
and rho has no such operation; the realisation here is therefore
different in kind, and is the technical content of this note.

We give the encoding, prove that its allocation traffic forms a tree
isomorphic to the parse tree of the source --- so that no channel is
shared by two allocation sites, which is the precise sense in which it is
not a distributed name server --- verify it against both counterexamples,
and prove it fully abstract for a bisimilarity whose observers are the
encodings of source contexts. That restriction is forced: an unrestricted
rho context separates the encodings of $P$ and $P \mid \nil$.
\end{abstract}

% =====================================================================
\section{Introduction}\label{sec:intro}
% =====================================================================

Process calculi customarily posit a countably infinite set of atomic
names and an operator $(\nu x)P$ producing one guaranteed fresh. An
implementation must discharge both, and in a distributed setting the
discharge is not routine.

The rho calculus \cite{MeredithRadestock05} moves the question inside the
theory. Names are quoted processes; the only binder is input; there is no
$\nu$. Whether this is a viable position is the question of whether the
$\pi$-calculus can be encoded into it without reintroducing what was
removed, and the original paper offered an encoding to that end. Its
design objective can be stated in one line:

\begin{quote}
\emph{no process in the image of the translation is a source of names for
any other.}
\end{quote}

\noindent
The price paid was that the translation compiles a whole term at a time
rather than acting clause by clause on a signature, because the namespace
discipline is a property of a term's structure.

Lybech \cite{Lybech22} showed that the encoding does not work.\footnote{%
The author discussed the matter with Lybech in early 2023 and suggested
that Abramsky's construction \cite{Abramsky93}, which had informed the
original design, also indicated the repair; the thread was not pursued.
The present note makes the argument in full.}
He gives two counterexamples, diagnoses their common cause exactly,
supplies a new encoding, and proves it valid against criteria adapted
from Gorla \cite{Gorla10} together with a criterion of \emph{parameter
independence} which is the right addition for parametrised translations.
None of that is in dispute here, and his diagnosis is adopted below
without amendment.

His encoding obtains freshness from a process
\[
  !N(x,z,v,s)
\]
reached on a fixed channel $v$ that every occurrence of $\nu$ in the term
must contact. That is a different implementation assumption --- a global
freshness oracle --- rather than a different source language, and under
that assumption his results stand. It is also an assumption he declines to
be comfortable with: his conclusion observes that for a distributed
system whose topology is not a star the translation does not represent
the source's structure, proposes instead that each replication instantiate
its own source of names, calls this closer to the intention of the 2005
encoding, and leaves it open, because such a scheme appears to need a
source of fresh names in order to build one.

\subsection*{The idea, in four lines}

There is no regress. Write $\Lft(\sigma) = \quo{(\out{\sigma}{\nil})}$
and $\Rgt(\sigma) = \quo{(\inp{y}{\sigma}{\nil})}$ for two structurally
distinct names built from $\sigma$. Then
\[
  \Frsh(\sigma;\,u,\,v).\,P
  \ \triangleq\
  \inp{u}{\sigma}{\inp{v}{\sigma}{P}}
  \ \mid\ \out{\sigma}{\out{\sigma}{\nil}}
  \ \mid\ \out{\sigma}{\inp{y}{\sigma}{\nil}}
\]
binds $u$ and $v$ to $\Lft(\sigma)$ and $\Rgt(\sigma)$ in two
communications, and it does so \emph{when $\sigma$ is a name received at
runtime}, because the two lifts mention $\sigma$ in their objects, which
substitution reaches, rather than under quotes, which it does not. A name
allocator is thus not a service to be contacted but a two-message
protocol that any node can run on a name it already holds. The whole
encoding is this construction applied at each node of the source term to
the address that node was given by its parent; the regress terminates at
a seed computed by the translation from the free names of the term and
consumed once.

Everything else in this note is an account of why $\Frsh$ has the shape it
has, why nothing simpler will do, and what it buys.

\subsection*{Contributions}

\begin{enumerate}[label=(\roman*),leftmargin=2.2em]

\item \textbf{Three facts about the rho calculus} (\S\ref{sec:facts}),
stated for the calculus rather than for any encoding. A reduction adds at
most one name to a term (Theorem~\ref{thm:onename}), and raises the
maximum quote depth of names in play by at most one
(Theorem~\ref{thm:onequote}); and substitution reaches the object of a
lift but not the interior of a quote
(Lemma~\ref{lem:transparency}). Between them these say that names are a
\emph{metered} resource: $n$ distinct new names cost $n$ rendezvous, and a
name $k$ quote levels above what a process holds costs $k$.

\item \textbf{A no-go result} (Corollary~\ref{cor:nogo}): no translation
of the 2005 shape can maintain distinctness across replication, since a
translation emits a finite term while replication requires unboundedly
many distinct names, and the duplicator reproduces quoted code verbatim
so that copies cannot be differentiated by syntax. Address arithmetic
must be communication.

\item \textbf{The discipline and its minimal realisation}
(\S\ref{sec:diagnosis}). Admissibility is a one-line syntactic side
condition on translation clauses which both broken clauses violate by
inspection; and any admissible allocator costs at least two
communications and two lifts mentioning its parameter, which $\Frsh$
attains (Proposition~\ref{prop:forced}).

\item \textbf{The encoding} (\S\ref{sec:encoding}), in which every
replication node carries a local address server derived from its own
address. Whole-term compilation survives: the allocator tree is static
and isomorphic to the parse tree; only the addresses are runtime values.

\item \textbf{A locality theorem} (\S\ref{sec:locality}) distinguishing
this from a distributed name server. The allocation traffic of $\enc{P}$
forms a tree isomorphic to the parse tree of $P$; no name is the subject
of allocations belonging to two distinct nodes; and there is in
particular no channel on which every allocation occurs. A server-based
encoding has, by contrast, a star.

\item \textbf{Full abstraction} (\S\ref{sec:fullabs}) for a
context-labelled bisimilarity whose observers are the encodings of source
contexts, together with a proof (\S\ref{sec:necessity}) that the
restriction is forced rather than convenient.

\end{enumerate}

\paragraph{Relation to Abramsky's construction.}
The architecture is his. The 2005 encoding already implemented one half
of it --- the splitting of a namespace across a multiplicative branch ---
and the invariant that half maintains is his, not a weaker one invented
for the occasion (\S\ref{sec:abramsky},
Proposition~\ref{prop:rightinv}). The other half, Copy's re-contraction
of a copied box's side ports, is what was missing. But it cannot be
transcribed: Copy renames the contents of a box at the metalevel, and rho
has no metalevel operation on quoted code --- this is not an oversight in
the calculus but the content of reflection, and the same fact on which
Lybech's separation theorem turns. Abramsky supplies the invariant and
the architecture; the realisation is necessarily different, and it is the
realisation that is at issue here.

\paragraph{Scope.}
The source language is the asynchronous, choice-free $\pi$-calculus with
input-guarded replication. Unguarded replication is not implementable ---
it runs away, and the 2005 encoding of it diverges unconditionally.
Lybech restricts to the same fragment for the same reason, and notes that
the restriction would not by itself have prevented either counterexample.
Cost accounting is out of scope.

% =====================================================================
\section{The rho calculus}\label{sec:rho}
% =====================================================================

\begin{definition}[Syntax]
Processes $P$ and names $x$ are given by mutual recursion:
\[
P \ ::=\ \nil \ \mid\ P_1 \mid P_2 \ \mid\ \out{x}{P} \ \mid\ \inp{y}{x}{P}
      \ \mid\ \drp{x}
\qquad\qquad
x \ ::=\ \quo{P}
\]
$\out{x}{P}$ is the \emph{lift}: it quotes $P$, creating the name
$\quo{P}$, and outputs it on $x$. $\inp{y}{x}{P}$ is input, binding $y$
in $P$; it is the only binder. $\drp{x}$ is the \emph{drop}: it strips
the quotes from $x$ and runs the process inside.
\end{definition}

\begin{definition}[Name equivalence, structural congruence, reduction]
$\nequiv$ is the least equivalence on names with
$\quo{P_1} \nequiv \quo{P_2}$ whenever $P_1 \equiv P_2$, and
$\quo{(\drp{x_1})} \nequiv x_2$ whenever $x_1 \nequiv x_2$;
$\equiv$ is the least congruence containing $\alpha$-equivalence and the
abelian monoid laws for $\mid$ with unit $\nil$, names being compared up
to $\nequiv$. Reduction is
\[
\frac{x_1 \nequiv x_2}
     {\inp{y}{x_1}{P_1} \ \mid\ \out{x_2}{P_2} \ \to\ P_1\{\quo{P_2}/y\}}
\]
closed under $\mid$ and $\equiv$.
\end{definition}

Substitution is capture-avoiding and replaces all names $\nequiv$-equal to
the substituted variable. Two clauses are used throughout:
\[
(\quo{P})\{x/y\} =
  \begin{cases} x & \text{if } \quo{P} \nequiv y\\ \quo{P} & \text{otherwise}\end{cases}
\qquad
(\drp{x})\{\quo{P}/y\} =
  \begin{cases} P & \text{if } x \nequiv y \\ \drp{x} & \text{otherwise}\end{cases}
\]

\begin{notation}
$\out{x}{\drp{z}}$ transmits the \emph{name} $z$ along $x$, since
$\quo{(\drp{z})} \nequiv z$. $\Dup(x) \triangleq
\inp{y}{x}{(\drp{y} \mid \out{x}{\drp{y}})}$ is the duplicator; it
satisfies $\out{x}{Q} \mid \Dup(x) \to Q \mid \out{x}{Q}$.
$\QD$ denotes Lybech's quote depth: $\QD(\quo{P}) = \QD(x)$ if
$P \equiv \drp{x}$, and $1 + \QD(P)$ otherwise, with $\QD(P)$ the maximum
of $\QD$ over the names occurring in $P$, or $0$ if there are none.
\end{notation}

\subsection{Why this notation}\label{sec:notation}

The 2005 papers wrote input as $x(y).P$ and quotation with corner quotes.
We use $\inp{y}{x}{P}$, $\quo{P}$, $\drp{x}$, for four reasons, the last
of which is the one that earns its place in a paper about a bug.

\begin{enumerate}[label=(\arabic*),leftmargin=2.2em]
\item It is ASCII; every term displayed here is a term of rholang, and
Appendix~\ref{app:rholang} runs.

\item It extends to joins without renegotiation. Simultaneous receipt is
$\inpp{y_1}{x_1}{y_2}{x_2}{P}$, with $\&$ joining within a group and $;$
separating alternatives. A prefix notation $x(y).P$ has no comparable
extension: the prefix is unary by construction. The encoding below does
not need joins; the generalisation to polyadic $\pi$
(\S\ref{sec:open}) does, and Remark~\ref{rem:joins} shows joins removing
a nondeterminism that the monadic version has to reason about.

\item It is the for-comprehension of Scala, Haskell's \texttt{do}, and
LINQ. A reader with no background in process calculi reads
$\inp{y}{x}{P}$ correctly at first sight: draw a $y$ from $x$, then do
$P$. This is not an analogy; RSpace, the store beneath the rho calculus
implementations, is a comprehension over a tuplespace.

\item $\quo{\cdot}$ and $\drp{\cdot}$ are visibly \emph{operators
occupying positions in a term}. Lemma~\ref{lem:transparency} is a
statement about which positions substitution can reach --- it reaches the
object of a lift, and not the interior of a quote --- and the whole
difference between the broken encoding and the repaired one is the
difference between $\out{c}{\out{n}{\nil}}$ and $\quo{(\out{n}{\nil})}$.
In corner quotes that difference is a change of typographic register.
\end{enumerate}

\begin{remark}
``rho'' abbreviates \textbf{r}eflective \textbf{h}igher \textbf{o}rder,
and the pun --- that it follows $\pi$ --- is the point of the name. We do
not write it with the letter $\rho$.
\end{remark}

% =====================================================================
\section{Three facts about name creation}\label{sec:facts}
% =====================================================================

This section is about the rho calculus and mentions no encoding. Its
content is that names are a metered resource.

\begin{definition}[Names in name position]
For a closed process $P$, let $\nms(P)$ be the set of names, up to
$\nequiv$, occurring in name position, defined by
\[
\begin{array}{r@{\ }c@{\ }l@{\qquad}r@{\ }c@{\ }l}
\nms(\nil) &=& \emptyset &
\nms(P_1 \mid P_2) &=& \nms(P_1) \cup \nms(P_2) \\
\nms(\out{x}{Q}) &=& \{x\} \cup \nms(Q) &
\nms(\inp{y}{x}{Q}) &=& \{x\} \cup (\nms(Q) \setminus \{y\}) \\
\nms(\drp{x}) &=& \{x\} .
\end{array}
\]
Note that $\quo{Q}$ is not in $\nms(\out{x}{Q})$: the lift's object is a
process, and the name it will become does not exist until the lift fires.
\end{definition}

\begin{lemma}\label{lem:nmscong}
$P \equiv Q$ implies $\nms(P) = \nms(Q)$, and
$\nms(P\{\quo{Q}/y\}) \subseteq (\nms(P) \setminus \{y\}) \cup \nms(Q) \cup \{\quo{Q}\}$.
\end{lemma}

\begin{proof}
The first by inspection of the axioms of $\equiv$. The second by
structural induction on $P$. The only cases that are not immediate are
$\drp{y}$, where $(\drp{y})\{\quo{Q}/y\} = Q$ and $\nms(Q)$ appears on
the right; and a name position occupied by $y$, which becomes $\quo{Q}$.
\end{proof}

\begin{theorem}[A reduction creates at most one name]\label{thm:onename}
If $P \to P'$ then $\nms(P') \subseteq \nms(P) \cup \{\quo{P_2}\}$, where
$P_2$ is the object of the lift consumed by the reduction. In particular
$|\nms(P')| \le |\nms(P)| + 1$.
\end{theorem}

\begin{proof}
By Lemma~\ref{lem:nmscong} we may work up to $\equiv$ and assume
$P = C[\,\inp{y}{x_1}{P_1} \mid \out{x_2}{P_2}\,]$ with
$P' = C[\,P_1\{\quo{P_2}/y\}\,]$. The context contributes the same names
to both sides. For the redex,
$\nms(\inp{y}{x_1}{P_1} \mid \out{x_2}{P_2}) =
 \{x_1,x_2\} \cup (\nms(P_1)\setminus\{y\}) \cup \nms(P_2)$, and by
Lemma~\ref{lem:nmscong}
$\nms(P_1\{\quo{P_2}/y\}) \subseteq (\nms(P_1)\setminus\{y\}) \cup \nms(P_2) \cup \{\quo{P_2}\}$.
\end{proof}

\begin{corollary}[Freshness costs communication]\label{cor:cost}
If $P \to^{k} P'$ then $|\nms(P') \setminus \nms(P)| \le k$. Hence a
process that exhibits $n$ pairwise non-$\nequiv$ names not among those it
starts with performs at least $n$ reductions.
\end{corollary}

\begin{theorem}[A reduction adds at most one quote level]\label{thm:onequote}
If $P \to P'$ then
$\max\{\QD(n) : n \in \nms(P')\} \le \max\{\QD(n) : n \in \nms(P)\} + 1$.
\end{theorem}

\begin{proof}
By Theorem~\ref{thm:onename} the only name in $\nms(P')$ possibly absent
from $\nms(P)$ is $\quo{P_2}$. If $P_2 \equiv \drp{x}$ then
$\quo{P_2} \nequiv x \in \nms(P)$ and there is nothing to prove.
Otherwise $\QD(\quo{P_2}) = 1 + \QD(P_2) = 1 + \max\{\QD(n) : n \in \nms(P_2)\}$,
and $\nms(P_2) \subseteq \nms(P)$.
\end{proof}

\begin{corollary}[One quote level per rendezvous]\label{cor:onequote}
Producing a name whose quote depth exceeds by $d$ the maximum quote depth
of the names a process holds requires at least $d$ reductions.
\end{corollary}

\begin{lemma}[Substitution transparency]\label{lem:transparency}
The object of a lift is a substitution-transparent position; a quote is
not:
\[
  (\out{c}{P})\{x/y\} \ =\ \out{c\{x/y\}}{(P\{x/y\})}
  \qquad\text{whereas}\qquad
  (\quo{P})\{x/y\} \in \{x,\ \quo{P}\}.
\]
Consequently a name may be derived from a runtime-bound $\sigma$ by the
process $\out{c}{\out{\sigma}{\nil}}$, which transmits $\Lft(\sigma)$ on
$c$; and it may not be derived by writing $\quo{(\out{\sigma}{\nil})}$ in
a term, since that expression is closed to substitution for $\sigma$.
\end{lemma}

\begin{proof}
Substitution is defined by structural recursion through process
constructors, and the lift's object is a process. Quotation is not a
process constructor but the name constructor, and its clause is the
two-case clause displayed above.
\end{proof}

\begin{remark}
Theorem~\ref{thm:onename} and Lemma~\ref{lem:transparency} pull in
opposite directions and together determine the shape of any name
discipline in the calculus. The first says that a name cannot be had for
free: it is paid for by a rendezvous. The second says where the payment
must be made: in the object of a lift, since that is the only position in
which a term mentioning a runtime parameter can still be reached by the
substitution that instantiates it.
\end{remark}

% =====================================================================
\section{The 2005 encoding and its objective}\label{sec:2005}
% =====================================================================

Fix an injection $\varphi$ from $\pi$ names to rho names; we suppress it.
The static quoting techniques generate from a name $\sigma$ the
\emph{increments}
\[
  \Lft(\sigma) = \quo{(\out{\sigma}{\nil})},
  \qquad
  \Rgt(\sigma) = \quo{(\inp{y}{\sigma}{\nil})},
\]
written $+\sigma$ and $\sigma{+}$ in the original. The translation
$\enc{-}_{n,p}$ takes two name parameters:
\[
\begin{array}{rcl}
\enc{\nil}_{n,p} &=& \nil \\
\enc{x(y).P}_{n,p} &=& \inp{y}{x}{\enc{P}_{n,p}} \\
\enc{\bar x \langle z\rangle}_{n,p} &=& \out{x}{\drp{z}} \\
\enc{P_1 \mid P_2}_{n,p} &=& \enc{P_1}_{\Lft(n),\Lft(p)} \mid \enc{P_2}_{\Rgt(n),\Rgt(p)} \\
\enc{(\nu z)P}_{n,p} &=& \inp{z}{p}{\enc{P}_{\Lft(n),\Lft(p)}} \mid \out{p}{\drp{n}} \\
\enc{!P}_{n,p} &=& \out{a}{\big(\inp{n}{b}{\inp{p}{c}{(\enc{P}_{n,p} \mid \Dup(a) \mid \cdots)}}\big)} \mid \Dup(a) \mid \cdots
\end{array}
\]
with $a = n\cdot p$, $b = \Rgt(n)$, $c = \Rgt(p)$, and top-level
parameters computed from $\fn(P)$ so that no name of $P$ is derivable
from them.

Two features are the design and should be recorded before the encoding is
criticised. First, no process in the image is a source of names for any
other: the top-level parameters are computed by the translation and never
requested. Second, the translation is parametrised by a position in the
term rather than by a head constructor, which is the price of the first
and the reason it is not a signature homomorphism.

The defect, in the form we shall need, is that the encoding carries an
invariant on $(n,p)$ which is nowhere written down.

% =====================================================================
\section{The counterexamples and the name server}\label{sec:lybech}
% =====================================================================

\begin{example}[Splitting]\label{ex:ce1}
For arbitrary $Q$,
$P_1' = !\big((\nu z)\bar u\langle z\rangle \mid Q\big)$ and
$P_2' = (\nu z)!\big(\bar u\langle z\rangle \mid Q\big)$
are separated in $\pi$ by a context that receives twice on $u$ and tests
the two names for equality; their images are not separated. The
replication clause rebinds $n$ and $p$ at each unfolding, but the clause
for $\mid$ has already frozen $\Lft(n)$ and $\Lft(p)$ inside quotes,
where the substitution cannot reach.
\end{example}

\begin{example}[Nesting]\label{ex:ce2}
$P_1 = !(\nu z)\bar u\langle z\rangle$ and
$P_2 = !(\nu q)(\nu z)\bar u\langle z\rangle$ are structurally
congruent; their images are separated, because the clause for $(\nu q)$
increments the parameters statically before the clause for $(\nu z)$ sees
them.
\end{example}

Lybech's diagnosis is exact and we adopt it: the property of being ``the
names most recently replicated'' is not preserved by the clause for
parallel composition, because the increments that clause writes are
static, placing the parameter under a quote, and substitution does not
recur under quotes.

\begin{remark}\label{rem:guarded}
Both examples use unguarded replication, which our fragment excludes;
Lybech records that the restriction would not have prevented either. We
restate them in the input-guarded fragment in \S\ref{sec:worked}.
\end{remark}

\paragraph{The name server.}
Lybech's encoding is $\enc{P}_{n,v} \mid\, !N(x,z,v,s)$, where
\[
  !N(x,z,v,s) \triangleq
  \Dup(x) \mid \out{x}{\big(\inp{a}{z}{\inp{r}{v}{(\Dup(x) \mid \out{r}{\drp{a}} \mid \out{z}{\out{a}{\nil}})}}\big)}
  \mid \out{z}{\drp{s}}
\]
holds a seed on $z$ and answers requests on the fixed channel $v$ with
successive left increments. The parameter $n$ is threaded forward and
never re-bound, which removes the second error; access to the name source
is by a channel that never changes, which removes the first. The
encoding is proved valid against compositionality, substitution
invariance, operational correspondence, observational correspondence,
divergence reflection and parameter independence.

Two observations about it, both of which we use.

\begin{enumerate}[label=(\arabic*),leftmargin=2.2em]
\item The implementation assumption is a global freshness oracle: one
process, reached on one fixed channel, is the source of every fresh name
in the term. The source language and the correctness results are not at
issue; the assumption is what differs from the objective of
\S\ref{sec:intro}, and \S\ref{sec:locality} makes the difference a
property of the image rather than a matter of emphasis.
\item The server builds its successor name by
$\out{z}{\out{a}{\nil}}$, where $a$ was \emph{received}. This computes
$\Lft(a)$ at runtime and works for the reason given in
Lemma~\ref{lem:transparency}. It is the manoeuvre the broken clauses
needed. Deployed once inside a server it looks like an implementation
detail; \S\ref{sec:diagnosis} makes it a discipline.
\end{enumerate}

% =====================================================================
\section{Abramsky's construction}\label{sec:abramsky}
% =====================================================================

Abramsky's computational interpretation of classical linear logic
\cite{Abramsky93} assigns \emph{proof expressions} $\Theta;\bar t$ to
sequent proofs and gives them a chemical-abstract-machine semantics. Four
features matter.

\paragraph{Names carry addresses.}
A bijection $\mathcal{N} \cong \mathbb{N}_0 \times \{l,r\}^*$ makes a name
a pair $\langle x_0, s\rangle$ of a root and a binary address; the
variants $t^l, t^r$ extend every address in $t$ by $l$ resp.\ $r$. The
invariant established by the proof-expression assignment and maintained by
every transition is: \emph{for distinct names $\langle x_0,s\rangle,
\langle y_0,t\rangle$ in $P$, $x_0 = y_0$ implies $s$ and $t$ are
incompatible in the prefix ordering.}

\paragraph{Contraction is a term former.}
$\displaystyle
\frac{\vdash \Theta; \Gamma,\ t{:}\,?A,\ u{:}\,?A}
     {\vdash \Theta; \Gamma,\ t \mathbin{@} u : \,?A}$
--- the joiner is a term, not a rewrite, and can be embedded and cut
against other terms.\footnote{Printed in \cite{Abramsky93} with $u{:}\,?B$;
comparison with the intuitionistic rule and with case (12) of the
realizability proof shows this to be a typo for $?A$.}

\paragraph{Copy.}
\[
  \bar x(P) \perp u \mathbin{@} v
  \ \longrightarrow\
  \bar x \perp (\bar x^{\,l} \mathbin{@} \bar x^{\,r}),\quad
  \bar x(P)^l \perp u,\quad
  \bar x(P)^r \perp v .
\]
Copying is demand-driven: the box duplicates only when a contraction asks
for two, and Copy is the only rule that increases the size of a proof
expression. Distinctness of the copies is by address extension, not
allocation: no fresh root is created, and roots are introduced only by the
name convention at Axiom, With and Of Course --- statically, finitely
often. The first clause is the rejoin: the box's side ports are cut
against a newly built contraction of the two copies' ports, so the
environment continues to see one port and \emph{the parent's namespace is
never re-rooted}.

\paragraph{What transfers and what does not.}
The address structure and the invariant transfer, and the 2005 encoding
already used them. The rejoin is what was missing, and it does not
transfer as a rule: $(\cdot)^l$ and $(\cdot)^r$ rewrite the interior of a
box, and by Lemma~\ref{lem:transparency} the rho calculus has no such
operation --- transmission of code without modification is what reflection
is. What can be carried over is the \emph{shape} of the rejoin: a port
shared by the copies, with the routing rather than the syntax doing the
distinguishing. That is what \S\ref{sec:encoding} builds.

\begin{center}
\begin{tabular}{@{}lll@{}}
\toprule
$\mathsf{PE}_2$ & 2005 encoding & \\
\midrule
$\mathcal{N} \cong \mathbb{N}_0 \times \{l,r\}^*$ & $\Lft, \Rgt$ & used \\
variants $t^l, t^r$ & the clause for $\mid$ & used \\
prefix-incompatibility invariant & --- & not stated \\
Contraction $t \mathbin{@} u$; Copy's first clause & --- & absent \\
demand-driven copying & eager replication (diverges) & absent \\
\bottomrule
\end{tabular}
\end{center}

% =====================================================================
\section{Diagnosis}\label{sec:diagnosis}
% =====================================================================

\begin{proposition}[The reconstructed invariant is the wrong one]
\label{prop:wronginv}
Let $I(n,p)$ hold when $(n,p)$ are the names most recently produced by
the enclosing replication context. No splitting clause preserves $I$: a
clause delegating to sub-translations at parameters derived from $(n,p)$
gives each sub-translation a pair which is a function of the most recent
output and not that output.
\end{proposition}

\begin{proposition}[Prefix-incompatibility is preserved by splitting]
\label{prop:rightinv}
Let $\mathrm{Inc}(S)$ hold when the addresses in $S$ are pairwise
incompatible in the prefix ordering. If $\mathrm{Inc}(S)$ and
$\sigma \in S$ then
$\mathrm{Inc}\big((S \setminus \{\sigma\}) \cup \{\sigma l, \sigma r\}\big)$.
\end{proposition}

\begin{proof}
$\sigma l, \sigma r$ are incompatible with each other; and any $\tau \in
S\setminus\{\sigma\}$ is incompatible with $\sigma$, hence with every
extension of $\sigma$.
\end{proof}

Taken together these say that the clause for parallel composition is not
the defect. Against the invariant the construction's ancestry supplies,
splitting is precisely the operation that maintains it. Against the
invariant the encoding was actually trying to maintain, nothing does.

\begin{corollary}[No static address assignment]\label{cor:nogo}
There is no translation with finitely many clauses, each emitting
finitely many quoted names, for which the name generated by a $\nu$ in the
$k$-th unfolding of a replication is written by the translation, for all
$k$.
\end{corollary}

\begin{proof}
The image of a fixed term is a finite term, so by
Theorem~\ref{thm:onename} the names it exhibits after $k$ reductions
number at most $|\nms(\enc{P})| + k$, and by
Corollary~\ref{cor:cost} $n$ distinct new names require $n$ rendezvous.
Hence they cannot all be written; they must be produced. The remaining
possibility is that a fixed piece of syntax is differentiated across
copies by substitution --- but the rho duplicator reproduces the quoted
body verbatim, and by Lemma~\ref{lem:transparency} substitution does not
reach beneath the quote in which that body sits.
\end{proof}

\begin{definition}[Admissibility]\label{def:discipline}
A translation clause is \emph{admissible} if every parameter bound at
runtime occurs in it only (a) in subject position of an input or a lift,
or (b) within the object of a lift; and never within a quote written by
the translation.
\end{definition}

\begin{proposition}
The clauses of \S\ref{sec:2005} for $\mid$ and for $\nu$ are
inadmissible; those for $\nil$, input and output are admissible.
\end{proposition}

\begin{proof}
By inspection: $\Lft(n) = \quo{(\out{n}{\nil})}$ places the parameter
under a written quote, and the two clauses use it. The remaining clauses
mention the parameters only by passing them on.
\end{proof}

\begin{proposition}[Admissible allocation is essentially $\Frsh$]
\label{prop:forced}
Let $A$ be an admissible process whose only free name is a parameter
$\sigma$ bound at runtime, and suppose $A$ has a reduction sequence after
which two names $u \not\nequiv v$, neither in $\nms(A)$, are bound in its
continuation. Then
\begin{enumerate}[label=(\roman*),leftmargin=2.2em,nosep]
\item $A$ performs at least two reductions; and
\item $A$ contains at least two lifts whose objects mention $\sigma$.
\end{enumerate}
$\Frsh$ attains both bounds.
\end{proposition}

\begin{proof}
(i) is Corollary~\ref{cor:cost}. For (ii), by
Theorem~\ref{thm:onename} each of the two names is $\quo{Q}$ for the
object $Q$ of some lift of $A$, and the two lifts are distinct since the
names differ. Suppose one such $Q$ does not mention $\sigma$. Then $Q$ is
closed, so $\quo{Q}$ is a constant of the translation, the same in every
instance of $A$; two instances of $A$ at distinct parameters would
allocate a common name, contradicting the requirement that the allocated
names not be among those already held once $A$ is placed in parallel with
another instance of itself --- which is exactly the situation inside a
replicated body. Hence both objects mention $\sigma$. $\Frsh$ has two
lifts and two reductions.
\end{proof}

\begin{remark}
Proposition~\ref{prop:forced} is why $\Frsh$ is displayed in
\S\ref{sec:intro} rather than derived at leisure. It is not a gadget
chosen for convenience: admissibility plus
Theorem~\ref{thm:onename} leave nothing else of that size, and the only
freedom remaining is the choice of the two templates, which
Lemma~\ref{lem:templates} constrains to be injective with disjoint
images.
\end{remark}

% =====================================================================
\section{The encoding}\label{sec:encoding}
% =====================================================================

\begin{definition}[Addresses]
An \emph{address} is a name $\theta(n_0)$ with
$\theta \in \{\Lft,\Rgt\}^*$ and $n_0$ the static seed
$\quo{\prod_{x \in \fn(P)} \out{x}{\nil}}$.
\end{definition}

We never write $\Lft$ or $\Rgt$ in the image of the translation; they are
metalanguage, used to \emph{name} addresses the encoding produces by
communication.

\begin{lemma}[Injectivity and disjointness]\label{lem:templates}
$\Lft$ and $\Rgt$ are injective up to $\nequiv$, their images are
disjoint, and $\QD(\Lft(\sigma)) = \QD(\Rgt(\sigma)) = 1 + \QD(\sigma)$.
Hence $\theta \mapsto \theta(n_0)$ is injective on $\{\Lft,\Rgt\}^*$.
\end{lemma}

\begin{proof}
$\quo{(\out{\sigma}{\nil})} \nequiv \quo{(\out{\tau}{\nil})}$ requires
$\out{\sigma}{\nil} \equiv \out{\tau}{\nil}$, hence $\sigma \nequiv \tau$;
similarly for $\Rgt$. A lift and an input are distinct constructors and
$\equiv$ does not relate them, giving disjointness. Neither template
yields a name of the form $\quo{(\drp{x})}$, so quote depth strictly
increases, which prevents an address from coinciding with a proper prefix
of itself; injectivity of $\theta \mapsto \theta(n_0)$ then follows by
induction on $|\theta|$.
\end{proof}

Lemma~\ref{lem:templates} realises Abramsky's bijection inside the rho
calculus with $n_0$ as the unique root, and resolves
prefix-incompatibility into two ingredients already present in the
literature: stratification by quote depth separates addresses of
different length, and template disjointness separates siblings.

\begin{construction}[The allocator]
\[
  \Frsh(\sigma;\, u,\, v).\,P
  \ \triangleq\
  \inp{u}{\sigma}{\inp{v}{\sigma}{P}}
  \ \mid\ \out{\sigma}{\out{\sigma}{\nil}}
  \ \mid\ \out{\sigma}{\inp{y}{\sigma}{\nil}}
\]
binding $u,v$ in $P$. It is admissible; after two reductions
$\{u,v\} = \{\Lft(\sigma),\Rgt(\sigma)\}$.
\end{construction}

That the assignment of $\Lft(\sigma)$ and $\Rgt(\sigma)$ to $u$ and $v$
is nondeterministic is deliberate. Only distinctness is ever required,
never identity; this is Lybech's criterion of parameter independence read
as a design principle instead of a proof obligation, and
Lemma~\ref{lem:addrind} discharges it once for the whole development.

\begin{definition}[The translation]\label{def:translation}
$\enc{P}$ is an \emph{address abstraction}, a function from a name to a
process; we write $\enc{P}_\sigma$ for its value.
\begin{align*}
\enc{\nil}_\sigma \ &=\ \nil \\[2pt]
\enc{\bar x\langle z\rangle}_\sigma \ &=\ \out{x}{\drp{z}} \\[2pt]
\enc{x(y).P}_\sigma \ &=\ \inp{y}{x}{\enc{P}_\sigma} \\[2pt]
\enc{P_1 \mid P_2}_\sigma \ &=\
  \inp{m}{\sigma}{\enc{P_1}_m}
  \ \mid\ \inp{m}{\sigma}{\enc{P_2}_m}
  \ \mid\ \out{\sigma}{\out{\sigma}{\nil}}
  \ \mid\ \out{\sigma}{\inp{y}{\sigma}{\nil}} \\[2pt]
\enc{(\nu z)P}_\sigma \ &=\ \Frsh(\sigma;\, z,\, m).\ \enc{P}_m \\[2pt]
\enc{!\,x(y).P}_\sigma \ &=\
  \Frsh(\sigma;\, a,\, b).\ \Frsh(b;\, s,\, c).\
  \big(\ \Dup(a) \ \mid\ \out{a}{B} \ \mid\ \out{s}{\drp{c}}\ \big)
\\
\text{where}\quad
B \ &=\ \inp{y}{x}{\Big(\Dup(a) \ \mid\
        \inp{c'}{s}{\ \Frsh(c';\, d,\, e).\ \big(\out{s}{\drp{e}} \mid \enc{P}_d\big)}\Big)}
\end{align*}
Top level: $\enc{P} \triangleq \enc{P}_{n_0}$.
\end{definition}

\paragraph{The replication clause.}
The node splits its address twice, obtaining a box channel $a$, a local
server port $s$ and a first seed $c$. The box $B$ is posted on $a$ and
duplicated in the self-reproducing manner, so
$\Dup(a) \mid \out{a}{B} \to B \mid \out{a}{B}$, whereupon $B$ blocks on
$x$: unfolding is input-guarded and there is no divergence. Each
unfolding reinstalls a duplicator, takes the current seed $c'$ off the
local port, splits it into a body address $d$ and a successor seed $e$,
returns $e$ to the port, and runs the body at $d$.

What is inside the quote is the point. By the time $B$ is posted, $a$,
$s$ and $x$ have been substituted, since they are bound by inputs whose
continuations include the lift, and by Lemma~\ref{lem:transparency}
substitution reaches a lift's object. Thereafter $B$ is copied verbatim
and needs no renaming, because it contains no addresses --- only the port
on which it asks for one. This is the shape of Copy's first clause: the
copies share the parent's port, and the routing does the distinguishing.

\begin{remark}[Locality]
The port $s$ is fixed and shared --- among the copies of \emph{one}
replication node, and derived from \emph{that node's} address. A term
with $k$ replication nodes has $k$ mutually unaware local servers.
\S\ref{sec:locality} makes this precise.
\end{remark}

\begin{remark}[Whole-term compilation survives]
Every process above is written by the translation. What became dynamic is
the \emph{value} of an address, not the structure computing it: the tree
of $\Frsh$ nodes is static and isomorphic to the parse tree.
\end{remark}

\begin{remark}[Joins]\label{rem:joins}
In a calculus with joins the two receipts of $\Frsh$ may be written
$\inpp{u}{\sigma}{v}{\sigma}{P}$, which consumes both messages in one
rendezvous and removes the nondeterminism entirely. We do not assume
joins in the theory --- the pure calculus does not have them --- but every
implementation does, and Appendix~\ref{app:rholang} uses the join form.
\end{remark}

% =====================================================================
\section{Locality: why this is not a distributed name server}
\label{sec:locality}
% =====================================================================

The obvious objection to \S\ref{sec:encoding} is that a name server per
replication node is still a name server. The objection is answered by
observing what the images actually communicate on.

\begin{definition}[Allocation site, name-traffic graph]
An \emph{allocation site} of an encoding $\enc{-}$ is an occurrence in
$\enc{P}$ of a subterm whose reduction produces a name not in
$\nms(\enc{P})$. The \emph{name-traffic graph} $G(\enc{P})$ has the
allocation sites as vertices, with an edge between two sites when a name
that is the subject of a communication at one is also the subject of a
communication at the other.
\end{definition}

\begin{theorem}[Locality]\label{thm:locality}
For every $\pi$ term $P$:
\begin{enumerate}[label=(\roman*),leftmargin=2.2em,nosep]
\item every subject occurring in $\enc{P}_{n_0}$ is either in
$\varphi(\fn(P))$, or is generated by a $\nu$ site, or is an address;
\item each address is the subject of communications belonging to exactly
one occurrence of a clause of Definition~\ref{def:translation};
\item the address that is the subject of an allocation is the one
allocated by the parent of the allocating node.
\end{enumerate}
\end{theorem}

\begin{proof}
(i) By induction on $P$: the clauses introduce no subjects beyond $x$,
$\sigma$, and names bound by their own inputs, and the latter are either
$\pi$ names received on $x$, or addresses received on $\sigma$.

(ii) Assign to each node of $P$ the address at which its clause is
invoked. The clauses for $\nil$, output and input allocate nothing and
pass their address to the continuation, so an address may be shared along
a prefix chain but is used as a subject only by the unique $\mid$, $\nu$
or $!$ node terminating that chain. Each of $\mid$, $\nu$ and $\Frsh$
uses its address $\sigma$ as a subject and allocates exactly
$\Lft(\sigma)$ and $\Rgt(\sigma)$, which by Lemma~\ref{lem:templates} lie
strictly below $\sigma$ in the address tree, so the subtrees allocated at
$\Lft(\sigma)$ and $\Rgt(\sigma)$ are disjoint. For $!$, the successive
seeds are $c_0 = c$ and $c_{k+1} = e_k$ where $d_k, e_k$ are the two
children of $c_k$; the body of the $k$-th unfolding therefore occupies the
subtree below $d_k$, and these are pairwise disjoint and disjoint from
every $c_j$. (Note that the successor seed must be a \emph{sibling} of
the body address: advancing by $\Lft(c_k)$ while the body ran at $c_k$
would make the next seed collide with the body's own first child.)

(iii) Immediate from the clauses: an allocating node's $\Frsh$ or split
has as subject the parameter it was passed.
\end{proof}

\begin{corollary}[No shared allocation channel]\label{cor:noserver}
$G(\enc{P})$ is isomorphic to the tree of allocating nodes of $P$. In
particular no name is the subject of allocations at two distinct nodes;
there is no name $v$ such that every allocation in $\enc{P}$ is a
communication on $v$; and the diameter of $G(\enc{P})$ grows with the
depth of $P$.
\end{corollary}

\begin{proof}
By Theorem~\ref{thm:locality}(ii) each address belongs to one node, and
by (iii) the only sharing is between a node and its parent, which is the
tree structure. A name $v$ as in the statement would be the subject of
allocations at every $\nu$ node, contradicting (ii) as soon as $P$ has
two.
\end{proof}

\begin{remark}
The same measurement applied to a server-based encoding gives a star: the
server channel $v$ is free in the image of every $\nu$, so $G$ has an
edge from each $\nu$ site to the server and its diameter is $2$
regardless of $P$. This is the difference between the two designs stated
as a property of the images rather than as a matter of emphasis, and it
is decidable by inspection. It is also, we think, the correct formal
content of Lybech's own observation that a star does not represent the
structure of a distributed source.
\end{remark}

\begin{corollary}[Freshness is confined to translation time]
No reduction of $\enc{P}$ consults a process that is not emitted by the
clause performing the reduction or by its parent; and the only name in
$\enc{P}$ not derived by reduction from another is $n_0$, which is
computed by the translation from $\fn(P)$.
\end{corollary}

\begin{remark}
Abramsky's construction is not free of a freshness assumption either: his
name convention stipulates a distinct root at each instance of Axiom,
With and Of Course. What it buys, and what the above inherits, is that
freshness is required only statically and finitely often. At runtime
there is no allocation from a supply --- only address extension,
performed locally, paid for at the rate fixed by
Theorem~\ref{thm:onename}.
\end{remark}

% =====================================================================
\section{The counterexamples, worked}\label{sec:worked}
% =====================================================================

We guard each replication with an input on a trigger channel $w$ and
supply triggers from the context.

\begin{example}[Example~\ref{ex:ce1}, guarded]
$P_1' = !\,w(t).\big((\nu z)\bar u\langle z\rangle \mid Q\big)$ and
$P_2' = (\nu z)\,!\,w(t).\big(\bar u\langle z\rangle \mid Q\big)$,
separated in $\pi$ by
$[\,\cdot\,] \mid \bar w \mid \bar w \mid u(n_1).u(n_2).(\bar n_1 \mid n_2.\bar x)$.

In $\enc{P_1'}$ the $k$-th unfolding runs its body at $d_k$, and the
body's parallel composition allocates $\Lft(d_k), \Rgt(d_k)$ by
communication on $d_k$, whose value differs with $k$ by
Theorem~\ref{thm:locality}(ii); the enclosed $\nu$ therefore emits a
different name at each unfolding. In $\enc{P_2'}$ the $\nu$ is outside
the box, its $\Frsh$ fires once, and every unfolding emits the same $z$.
The offending static increment has become a communication, and the value
it computes tracks the unfolding.
\end{example}

\begin{example}[Example~\ref{ex:ce2}, guarded]
$P_1 = !\,w(t).(\nu z)\bar u\langle z\rangle$ and
$P_2 = !\,w(t).(\nu q)(\nu z)\bar u\langle z\rangle$, with
$P_1 \equiv P_2$. Unfolding $k$ of $\enc{P_1}$ runs
$\Frsh(d_k;z,m).\out{u}{\drp{z}}$ and emits a child of $d_k$; unfolding
$k$ of $\enc{P_2}$ runs $\Frsh(d_k;q,m).\Frsh(m;z,m').\out{u}{\drp{z}}$
and emits a grandchild. In both the emitted names are pairwise distinct
across $k$, since the $d_k$ are, and in neither is any emitted name
observable except as a $\pi$ name; the two are identified, as they must
be.
\end{example}

% =====================================================================
\section{Full abstraction relative to encoded contexts}\label{sec:fullabs}
% =====================================================================

\subsection{Context-labelled bisimilarity}

\begin{definition}
For a calculus with reduction $\to$ and contexts $\Ctx$, write
$P \xrightarrow{\,C\,} P'$ when $C[P] \to P'$, and
$P \xLongrightarrow{\,C\,} P'$ when $C[P] \to^* P'$. For
$\Obs \subseteq \Ctx$, a symmetric $\mathcal{R}$ is a
$\Obs$-bisimulation if $P \mathrel{\mathcal{R}} Q$ and
$P \xrightarrow{\,C\,} P'$ with $C \in \Obs$ imply
$Q \xLongrightarrow{\,C\,} Q'$ with $P' \mathrel{\mathcal{R}} Q'$;
$\bisim_\Obs$ is the largest such.
\end{definition}

On the source $\Obs = \Ctx_\pi$, written $\bisim_\pi$; on the target
$\Obs = \enc{\Ctx_\pi}$, the encodings of source contexts.

\begin{remark}[On minimality]
The labels are enabling contexts, not \emph{minimal} enabling contexts.
Minimality --- relative pushouts, and their refinement to a subcategory
of admissible observers --- is the right long-run discipline; it is
developed in the companion paper \cite{BIHOL} and taken up in
\S\ref{sec:future}, and nothing below needs it. Without minimality the
label set is larger and $\bisim_\Obs$ correspondingly finer, which makes
completeness harder rather than easier.
\end{remark}

\begin{definition}
$\enc{-}$ extends to contexts with $\enc{[\,\cdot\,]}_\sigma$ a hole
expecting an address; plugging is instantiation of the address
abstraction.
\end{definition}

\begin{lemma}[Compositionality]\label{lem:compositional}
$\enc{C[P]} = \enc{C}[\enc{P}]$, on the nose.
\end{lemma}

\begin{proof}
Induction on $C$: each clause of Definition~\ref{def:translation}
mentions its sub-translations only under an address abstraction, and the
hole occurs in exactly one.
\end{proof}

\subsection{Name usage and relabelling}

The next three results are the apparatus the bisimulation argument rests
on. The first is the syntactic induction that makes the informal claim
``encoded contexts do not inspect names'' precise.

\begin{lemma}[Name usage]\label{lem:usage}
Let $y$ be a variable bound in $\enc{P}_\sigma$ by an input arising from
a $\pi$ input, or by a $\Frsh$ in the position of a $\nu$-bound name.
Then every occurrence of $y$ in $\enc{P}_\sigma$ is either the subject of
an input or of a lift, or is the $z$ in a subterm $\out{x}{\drp{z}}$. In
particular no occurrence of $y$ lies inside a quote, and no subterm
$\drp{y}$ occurs except under a lift as above.
\end{lemma}

\begin{proof}
Induction on $P$. For $\bar x\langle z\rangle$ the image is
$\out{x}{\drp{z}}$, in which $x$ is a subject and $z$ is as described.
For $x(y).P$ the image is $\inp{y}{x}{\enc{P}_\sigma}$ and the claim for
$y$ follows from the hypothesis for $P$. For $(\nu z)P$ the image is
$\Frsh(\sigma; z, m).\enc{P}_m$, and $z$ occurs only in
$\enc{P}_m$. For $\mid$ and for $!$ the bound variables introduced are
addresses, not $\pi$ names, and are used only as subjects or in lift
objects by construction (admissibility, Definition~\ref{def:discipline}).
\end{proof}

\begin{definition}[Address relabelling]
For addresses $\sigma,\tau$, let $\beta_{\sigma\to\tau}$ be the map on
names sending $\theta(\sigma) \mapsto \theta(\tau)$ for
$\theta \in \{\Lft,\Rgt\}^*$ and fixing all other names. It is a
bijection by Lemma~\ref{lem:templates}, and extends to processes by
metalevel replacement.
\end{definition}

\begin{remark}
$\beta_{\sigma\to\tau}$ is \emph{not} a rho substitution: it must descend
into the interiors of quotes, and by Lemma~\ref{lem:transparency} no
substitution does. It is available to the metatheory and to no rho
context. This is a compact statement of why the encoding is safe from
encoded observers, and \S\ref{sec:necessity} is a statement of why it is
not safe from arbitrary ones.
\end{remark}

\begin{lemma}[Equivariance]\label{lem:equivariance}
$\beta_{\sigma\to\tau}(\enc{P}_\sigma) = \enc{P}_\tau$ whenever $\sigma,
\tau$ satisfy the distinctness condition of
Theorem~\ref{thm:locality}.
\end{lemma}

\begin{proof}
Induction on $P$, using that every occurrence of an address in
$\enc{P}_\sigma$ is either $\sigma$ itself or a bound variable later
instantiated to $\theta(\sigma)$, and that
$\beta$ commutes with $\Lft$ and $\Rgt$ by construction.
\end{proof}

\begin{lemma}[Address independence]\label{lem:addrind}
$\enc{P}_\sigma \bisim_{\enc{\Ctx_\pi}} \enc{P}_\tau$ for
$\sigma,\tau$ as above.
\end{lemma}

\begin{proof}
Take $\mathcal{R} = \{(T, \beta_{\sigma\to\tau}(T))\}$ for $T$ reachable
from $\enc{P}_\sigma$. Lemma~\ref{lem:equivariance} gives the base case.
For the step, $\beta$ is a bijection on names, so it preserves and
reflects $\nequiv$ and hence the applicability of the reduction rule;
and for the labels, Lemma~\ref{lem:usage} shows that an encoded context
uses a name it has received only as a subject or as the object of a lift
of the form $\out{x}{\drp{z}}$, so the only discrimination it can make
between two received names is an equality test by attempted
synchronisation, which a bijection preserves and reflects. In particular
no encoded context applies $\drp{\cdot}$ to a received name, which is the
only operation that could observe a name's structure.
\end{proof}

\begin{corollary}[The allocator race is harmless]\label{cor:race}
The two reductions of a $\Frsh$ commute up to
$\bisim_{\enc{\Ctx_\pi}}$: the two outcomes differ by
$\beta$ exchanging $\Lft(\sigma)$ and $\Rgt(\sigma)$ and their subtrees.
\end{corollary}

\subsection{Administrative reductions}

\begin{lemma}[Administrative reductions terminate]\label{lem:admin}
Call a reduction \emph{administrative} if its subject is an address.
Every maximal administrative sequence from $\enc{P}_\sigma$ is finite, and
no administrative reduction is observable on $\varphi(\fn(P))$.
\end{lemma}

\begin{proof}
Each $\mid$ node contributes two, each $\nu$ two, each $!$ four plus, per
unfolding, one on $a$ and three on $s$ and $c_k$; unfoldings are guarded
by an input on the source channel, so the number of administrative steps
enabled is bounded by a function of the number of source communications
performed. Unobservability is Theorem~\ref{thm:locality}(i) together with
the choice of $n_0$.
\end{proof}

\begin{corollary}[Divergence reflection]
$\enc{P} \to^{\omega}$ implies $P \to^{\omega}$.
\end{corollary}

\subsection{Operational correspondence and full abstraction}

\begin{proposition}[Completeness]\label{prop:complete}
If $P \to P'$ then $\enc{P}_\sigma \to^{+} T$ with
$T \bisim_{\enc{\Ctx_\pi}} \enc{P'}_{\sigma'}$ for some address
$\sigma'$.
\end{proposition}

\begin{proposition}[Soundness]\label{prop:sound}
If $\enc{P}_\sigma \to^* T$ then $P \to^* P'$ for some $P'$ with
$T \bisim_{\enc{\Ctx_\pi}} \enc{P'}_{\sigma'}$.
\end{proposition}

\begin{theorem}[Full abstraction]\label{thm:fullabs}
$P \bisim_\pi Q \iff \enc{P} \bisim_{\enc{\Ctx_\pi}} \enc{Q}$.
\end{theorem}

\begin{proposition}\label{prop:criteria}
Theorem~\ref{thm:fullabs} implies Lybech's criteria 1--6, with $\simeq$
taken to be $\bisim_{\enc{\Ctx_\pi}}$ and the name derivation function
$\sigma \mapsto \{\Lft(\sigma),\Rgt(\sigma)\}$.
\end{proposition}

The proofs of Propositions~\ref{prop:complete}, \ref{prop:sound},
\ref{prop:criteria} and Theorem~\ref{thm:fullabs} are given in
Appendix~\ref{app:proofs}. The apparatus they use is
Lemmas~\ref{lem:compositional}, \ref{lem:usage}, \ref{lem:equivariance},
\ref{lem:addrind} and \ref{lem:admin} and no more; in particular the
nondeterminism of $\Frsh$ enters only through
Corollary~\ref{cor:race}.

\subsection{The restriction is forced}\label{sec:necessity}

\begin{proposition}\label{prop:necessity}
There are $\pi$ processes $P \equiv Q$ and a rho context $C \notin
\enc{\Ctx_\pi}$ with $C[\enc{P}] \not\bisim C[\enc{Q}]$.
\end{proposition}

\begin{proof}
Take $P = \bar u\langle v\rangle$ and
$Q = \bar u\langle v\rangle \mid \nil$. Then
$\enc{P}_{n_0} = \out{u}{\drp{v}}$, with no barb on $n_0$, whereas
\[
\enc{Q}_{n_0} =
  \inp{m}{n_0}{\out{u}{\drp{v}}} \ \mid\ \inp{m}{n_0}{\nil}
  \ \mid\ \out{n_0}{\out{n_0}{\nil}} \ \mid\ \out{n_0}{\inp{y}{n_0}{\nil}}
\]
has one. Since $n_0$ is computed by the translation from $\fn(P)$, an
arbitrary rho context can name it: take
$C = [\,\cdot\,] \mid \inp{k}{n_0}{\out{\mathit{succ}}{\nil}}$. No
encoded context can, since $n_0 \notin \varphi(\fn(P))$.
\end{proof}

\begin{remark}[The sharper form]
A deeper version exploits reflection rather than scaffolding: an
arbitrary rho context may apply $\drp{\cdot}$ to a name received on $u$,
and since a $\nu$-generated name here is $\Lft(\sigma)$ or $\Rgt(\sigma)$
for an internal $\sigma$, dropping it runs $\out{\sigma}{\nil}$ or
$\inp{y}{\sigma}{\nil}$ and exposes a barb on an internal address. By
Lemma~\ref{lem:usage} no encoded context does this. This is not a
weakness of the present encoding: by Lybech's separation theorem the rho
calculus manufactures new free --- hence observable and substitutable ---
names at runtime, and no encoding of $\pi$ into rho can be fully abstract
against unrestricted rho observers. The restriction to $\enc{\Ctx_\pi}$
is the shape of the theorem rather than a weakening of it.
\end{remark}

\begin{remark}
Lybech's $\bisim^{\mathcal{N}}$ restricts observation by stipulating a set
of names; the restriction here is derived from the source. The two agree
for $\pi$ and rho, both being name-passing calculi for which barbs are an
adequate observation. They do not agree in general: a name-set
presentation assumes that what one can observe of a term is the set of
names it can communicate on, which is false for, e.g., the ambient
calculus.
\end{remark}

% =====================================================================
\section{Discussion}\label{sec:discussion}
% =====================================================================

\begin{center}
\begin{tabular}{@{}p{0.32\textwidth}p{0.28\textwidth}p{0.28\textwidth}@{}}
\toprule
 & name server & this note \\
\midrule
implementation assumption & a global freshness oracle on a fixed channel
  & no runtime name source; a static seed \\
name-traffic graph & star, diameter $2$ & tree $\cong$ parse tree \\
translation & clause-by-clause & whole term \\
where freshness is required & at every $\nu$, at runtime & at translation
  time only \\
\bottomrule
\end{tabular}
\end{center}

The table compares design objectives and the implementation assumptions
that follow from them, not degrees of correctness: both encodings are
correct under their own assumptions, for the same source language, and
Lybech's is proved in more detail than this one currently is.

\paragraph{Demand-driven copying is a dividend.}
Abramsky's Copy fires only when a contraction asks for a second copy.
Restricting to input-guarded replication, which both encodings do to
avoid divergence, is the process-calculus shadow of that discipline; the
fix for divergence and the fix for freshness come from the same place.

\paragraph{Linearity was doing work we must now do by hand.}
In $\mathsf{PE}_2$ the address invariant is protected by linearity --- each
name occurs exactly twice --- and by acyclicity, which gives deadlock
freedom. The rho calculus has neither. What was structural there is a
theorem here, and its proof, Lemmas~\ref{lem:templates} and
Theorem~\ref{thm:locality}, decomposes along exactly Abramsky's two axes:
depth and siblings.

% =====================================================================
\section{Related work}\label{sec:related}
% =====================================================================

Abramsky's programme was carried forward by Bellin and Scott
\cite{BellinScott94} and tightened into a session-typed discipline by
Caires and Pfenning \cite{CairesPfenning10} and Wadler \cite{Wadler14}.
Kokke, Montesi and Peressotti \cite{HCP} observe that the separation of
environments internalised by $\otimes$ corresponds to parallel process
behaviour and extend linear logic to hyperenvironments accordingly: the
same move as the present one, made from the side of the logic.

On encodability, Gorla \cite{Gorla10} supplies the criteria Lybech
adapts; Gorla and Nestmann \cite{GorlaNestmann14} argue against full
abstraction as a criterion; van Glabbeek \cite{Glabbeek18} derives
validity from a semantic equivalence rather than a list, which is closer
in spirit to \S\ref{sec:fullabs}, though he does not treat parametrised
translations. Sangiorgi's encoding of HO$\pi$ into $\pi$
\cite{Sangiorgi93} is the result Lybech's separation theorem shows does
not extend to reflection; the relabelling apparatus of
\S\ref{sec:fullabs} is the standard treatment of distinctions and indexed
bisimulation \cite{SangiorgiWalker01} and is cited rather than
reconstructed. Bendixen, Bojesen, H\"uttel and Lybech \cite{Psi}
instantiate the rho calculus as a higher-order $\Psi$-calculus. Leifer
and Milner \cite{LeiferMilner00} supply the derived-transition machinery
a minimal-label refinement would use; \cite{BIHOL} refines it to a
subcategory of admissible observers, which is the form \S\ref{sec:future}
requires.

% =====================================================================
\section{Open problems}\label{sec:open}
% =====================================================================

\begin{enumerate}[label=\arabic*.,leftmargin=2.2em]
\item \textbf{Minimal labels.} Replace enabling contexts by minimal
ones, taken relative to $\enc{\Ctx_\pi}$ rather than to all of rho, and
show Theorem~\ref{thm:fullabs} survives with a coarser relation. See
\S\ref{sec:future}(ii) for the shape of the transfer and \cite{BIHOL}
for the apparatus.
\item \textbf{A converse to Theorem~\ref{thm:onename}.} The theorem
bounds name creation from above. A matching lower bound for a class of
allocation protocols would turn Proposition~\ref{prop:forced} into a
genuine optimality result rather than a bound attained.
\item \textbf{Polyadic $\pi$ and joins.} Encode polyadic synchronisation
using the join form directly, connecting with Carbone and Maffeis
\cite{CarboneMaffeis03}.
\item \textbf{Choice.} The source fragment is choice-free.
\item \textbf{Cost.} The $\Frsh$ tree is a metering surface: each source
$\nu$ costs two communications, each unfolding four.
Theorem~\ref{thm:onename} says the rate is optimal; what a whole term
costs is not treated here.
\item \textbf{Mechanisation.} Definition~\ref{def:discipline} is a
syntactic side condition and should be machine checked; a translation
satisfying it cannot have a bug of the kind found in 2022.
\end{enumerate}

% =====================================================================
\section{Future work: contexts as morphism data}\label{sec:future}
% =====================================================================

The equivalence of \S\ref{sec:fullabs} is indexed by a class of
contexts, and that class was obtained by applying the translation to the
contexts of the source. Two maps are therefore in play and only one of
them has been named: $\enc{-}$ on terms, and $\enc{-}$ on contexts. The
present note gets away with conflating them because
Lemma~\ref{lem:compositional} holds on the nose --- the translation of a
plugged context is the plugged translation of the context --- but the
conflation is an accident of this encoding rather than a general fact,
and it is worth saying what the general statement should be.

The observation is that in a context-labelled transition system the
labels \emph{are} contexts. A map of theories that is asked to preserve
context-labelled bisimilarity is therefore being asked to do something a
map on terms cannot by itself do: it must say, of a source context, which
target context is to serve as its label. Two translations agreeing on
every term may realise a given source context differently and preserve
different relations. What is missing from the received notion of a
behaviour-preserving translation is thus not a hypothesis but a
component. The shape we expect is a pair $(F,\Phi)$ --- $F$ on terms,
$\Phi$ on contexts --- subject to an equivariance condition
\[
  F(c[t_1,\dots,t_n]) \ \approx\ \Phi(c)[F t_1,\dots,F t_n]
\]
and to preservation of transitions along $\Phi$; the class of admissible
observers is then not a parameter chosen by the analyst but the image of
$\Phi$, and the restriction to $\enc{\Ctx_\pi}$ that
\S\ref{sec:necessity} shows to be forced becomes the instance
$\Obs = \operatorname{im}\Phi$ of a definition rather than a stipulation
attached to a theorem.

Three things in this note are evidence about that definition, and we
record them as such.

\begin{enumerate}[label=(\roman*),leftmargin=2.2em]

\item \textbf{The strict case is not vacuous.} The equivariance
condition above is stated up to $\approx$ because the $\nu$ equations of
a source calculus are in general only satisfied up to bisimilarity in a
target. For the encoding of Definition~\ref{def:translation} it holds on
the nose (Lemma~\ref{lem:compositional}). A nontrivial encoding with
$\Phi$ strict is worth having, since it shows that the lax formulation is
demanded by particular examples rather than by the shape of the
definition.

\item \textbf{Minimality is the missing discipline, and it is
$\Phi$-relative.} Open problem~1 asks to replace enabling contexts by
minimal ones. The right notion of minimality is not minimality in the
target --- the minimal rho context enabling a transition out of $\enc{P}$
is typically smaller than any $\enc{C}$, because $\enc{C}$ carries
bookkeeping that minimisation strips away --- but minimality within
$\operatorname{im}\Phi$. Relative pushouts taken in that subcategory are
the construction this calls for, and the refinement of the
Leifer--Milner apparatus that supplies them is developed in a companion
paper \cite{BIHOL}. The transfer is asymmetric and worth stating
precisely: a larger label set gives a finer relation, so the forward
half of Theorem~\ref{thm:fullabs} descends to the minimal-label relation
for free, while the backward half does not. What is owed is a lemma
saying that relative minimality loses no separating power in the image of
$\Phi$; it is a statement about relative pushouts and belongs with them.

\item \textbf{Address relabelling is a symmetry of
$\operatorname{im}\Phi$ and not of the target.} The map
$\beta_{\sigma\to\tau}$ descends into quotes, so it is not a rho
substitution and no rho context performs it; it is nonetheless a
bisimulation for encoded observers (Lemma~\ref{lem:addrind}). A general
account should predict such symmetries from $\Phi$ rather than exhibit
them per encoding, which is the same demand as asking that the generated
spatial-behavioural logics be functorial on the morphism class.

\end{enumerate}

We intend to develop this in the companion paper, where the encoding
given here is the worked instance, and to return with the minimal-label
form of Theorem~\ref{thm:fullabs}.

\section*{Acknowledgements}

The author thanks Stian Lybech, whose counterexamples are correct and
whose diagnosis of their common cause is the reason the present repair
can be stated briefly. Claude (Anthropic) assisted with source review,
drafting, and proof-checking.

\begin{thebibliography}{99}

\bibitem{Abramsky93} S. Abramsky.
\emph{Computational interpretations of linear logic.}
Theoretical Computer Science 111(1--2):3--57, 1993.

\bibitem{BellinScott94} G. Bellin and P. J. Scott.
\emph{On the $\pi$-calculus and linear logic.}
Theoretical Computer Science 135(1):11--65, 1994.

\bibitem{BIHOL} C.\,B. Wells, M. Stay and L.\,G. Meredith.
\emph{Behavior in higher-order languages: admissible observers, relative
pushouts, and the encoding of extended theories.}
In preparation, F1R3FLY.io, 2026.

\bibitem{CairesPfenning10} L. Caires and F. Pfenning.
\emph{Session types as intuitionistic linear propositions.}
CONCUR 2010, LNCS 6269, 222--236.

\bibitem{CarboneMaffeis03} M. Carbone and S. Maffeis.
\emph{On the expressive power of polyadic synchronisation in $\pi$-calculus.}
Nordic Journal of Computing 10(2):70--98, 2003.

\bibitem{Glabbeek18} R. van Glabbeek.
\emph{A theory of encodings and expressiveness.}
FoSSaCS 2018, 183--202.

\bibitem{Gorla10} D. Gorla.
\emph{Towards a unified approach to encodability and separation results
for process calculi.}
Information and Computation 208(9):1031--1053, 2010.

\bibitem{GorlaNestmann14} D. Gorla and U. Nestmann.
\emph{Full abstraction for expressiveness: history, myths and facts.}
Mathematical Structures in Computer Science 26:639--654, 2014.

\bibitem{HCP} W. Kokke, F. Montesi and M. Peressotti.
\emph{Better late than never: a fully abstract semantics for classical
processes.}
Proc. ACM Program. Lang. 3(POPL), 2019.

\bibitem{LeiferMilner00} J. J. Leifer and R. Milner.
\emph{Deriving bisimulation congruences for reactive systems.}
CONCUR 2000, LNCS 1877, 243--258.

\bibitem{Lybech22} S. Lybech.
\emph{Encodability and separation for a reflective higher-order calculus.}
EXPRESS/SOS 2022, EPTCS 368, 95--112.

\bibitem{MeredithRadestock05} L. G. Meredith and M. Radestock.
\emph{A reflective higher-order calculus.}
Electronic Notes in Theoretical Computer Science 141(5):49--67, 2005.

\bibitem{Psi} A. R. Bendixen, B. B. Bojesen, H. H\"uttel and S. Lybech.
\emph{A generic type system for higher-order $\Psi$-calculi.}
EXPRESS/SOS 2022, EPTCS 368, 43--59.

\bibitem{Sangiorgi93} D. Sangiorgi.
\emph{From $\pi$-calculus to higher-order $\pi$-calculus --- and back.}
TAPSOFT 1993, LNCS 668, 151--166.

\bibitem{SangiorgiWalker01} D. Sangiorgi and D. Walker.
\emph{The $\pi$-calculus: a theory of mobile processes.}
Cambridge University Press, 2001.

\bibitem{Wadler14} P. Wadler.
\emph{Propositions as sessions.}
Journal of Functional Programming 24(2--3):384--418, 2014.

\end{thebibliography}

\appendix

% =====================================================================
\section{Status of results}\label{app:status}
% =====================================================================

\begin{center}
\begin{tabular}{@{}llp{0.40\textwidth}@{}}
\toprule
Result & Status & Note \\
\midrule
Thm.~\ref{thm:onename}, Cor.~\ref{cor:cost} & proved & \\
Thm.~\ref{thm:onequote}, Cor.~\ref{cor:onequote} & proved & \\
Lem.~\ref{lem:transparency} & proved & \\
Prop.~\ref{prop:wronginv}, \ref{prop:rightinv} & proved & elementary \\
Cor.~\ref{cor:nogo} & proved & from \S\ref{sec:facts} \\
Prop.~\ref{prop:forced} & proved & bound attained, not optimality; see open problem 2 \\
Lem.~\ref{lem:templates} & proved & uses Lybech's stratification lemma \\
Thm.~\ref{thm:locality}, Cor.~\ref{cor:noserver} & proved & \\
Lem.~\ref{lem:compositional} & proved & \\
Lem.~\ref{lem:usage} & proved & syntactic induction \\
Lem.~\ref{lem:equivariance}, \ref{lem:addrind}, Cor.~\ref{cor:race} & proved & \\
Lem.~\ref{lem:admin} & proved & \\
Prop.~\ref{prop:necessity} & proved & \\
Prop.~\ref{prop:subst} & proved & App.~\ref{app:proofs}; uses Def.~\ref{def:discipline} \\
Lem.~\ref{lem:commute}, \ref{lem:seed} & proved & App.~\ref{app:proofs} \\
Lem.~\ref{lem:exposure} & proved & App.~\ref{app:proofs}; the structural lemma \\
Prop.~\ref{prop:complete}, \ref{prop:sound} & proved & App.~\ref{app:proofs} \\
Thm.~\ref{thm:fullabs}, Prop.~\ref{prop:criteria} & proved & App.~\ref{app:proofs} \\
\bottomrule
\end{tabular}
\end{center}

% =====================================================================
\section{Operational correspondence and full abstraction}\label{app:proofs}
% =====================================================================

This appendix supplies the proofs of Propositions~\ref{prop:complete},
\ref{prop:sound} and \ref{prop:criteria} and of
Theorem~\ref{thm:fullabs}. The apparatus is that of \S\ref{sec:fullabs}
--- Lemmas~\ref{lem:compositional}, \ref{lem:usage},
\ref{lem:equivariance}, \ref{lem:addrind}, \ref{lem:admin} and
Theorem~\ref{thm:locality} --- together with one device introduced here:
a relation $\real$ between rho processes and source configurations,
which says that a target term is \emph{some} state of the machine
realising a given source term, not necessarily a state in which the
translation's own bookkeeping has finished. It is that latitude that the
replication clause needs, since a source communication may fire while a
box copy is complete and the seed exchange is not.

\subsection{Address environments}

The translation of Definition~\ref{def:translation} sends a $\pi$ name to
itself; the top-level map $\varphi$ fixes an injection of source names
into rho names, and $\nu$-bound names are instantiated at runtime by
$\Frsh$. To speak of a term in mid-computation we must record which
addresses the $\nu$'s have produced so far.

\begin{definition}[Address environment]\label{def:aenv}
An \emph{address environment} $\rho$ is a finite partial injection from
$\pi$ names to addresses such that no address in its image is a prefix of
another and none is a prefix of, or prefixed by, an address in use as a
translation parameter. Write $\hat\rho$ for the total map sending
$x \mapsto \rho(x)$ when $x \in \dom\rho$ and $x \mapsto \varphi(x)$
otherwise, and
\[
  \enc{P}^\rho_\sigma \ \triangleq\ \enc{P}_\sigma\,\hat\rho
\]
for the result of replacing every free occurrence of a $\pi$ name $x$ in
$\enc{P}_\sigma$ by $\hat\rho(x)$.
\end{definition}

\begin{lemma}\label{lem:aenv-subst}
$\enc{P}_\sigma\,\hat\rho$ is a rho substitution: every free occurrence of
a $\pi$ name in $\enc{P}_\sigma$ lies in a position that substitution
reaches.
\end{lemma}

\begin{proof}
By Lemma~\ref{lem:usage} such an occurrence is the subject of an input,
the subject of a lift, or the $z$ of a subterm $\out{x}{\drp{z}}$; none of
these is inside a quote. The only sub-image of the translation that lies
inside a quote is the box $B$ of the replication clause, which occurs as
the object of the lift $\out{a}{B}$, and by
Lemma~\ref{lem:transparency} substitution reaches a lift's object.
\end{proof}

\begin{convention}
Throughout, $\sigma$ is an address satisfying the distinctness condition
of Theorem~\ref{thm:locality} with respect to $\rho$, and we write
$\enc{P}^\rho_\sigma$ only when $\fn(P) \subseteq \dom\rho \cup \dom\varphi$.
Configurations are triples $(P,\rho,\sigma)$ so constrained.
\end{convention}

\subsection{B.1 Substitution}

\begin{proposition}[Substitution]\label{prop:subst}
Let $y$ be free in $P$ and let $n$ be a rho name not occurring in
$\enc{P}_\sigma$ as an address. Then
\[
  \enc{P}^\rho_\sigma\{n/y\} \ = \ \enc{P}^{\rho[y\mapsto n]}_\sigma .
\]
In particular $\enc{P\{z/y\}}^\rho_\sigma = \enc{P}^\rho_\sigma\{\hat\rho(z)/y\}$.
\end{proposition}

\begin{proof}
Induction on $P$, using Lemma~\ref{lem:aenv-subst} to know that no
occurrence of $y$ is shielded by a quote.

$P = \nil$: both sides are $\nil$.

$P = \bar x\langle z\rangle$: the image is $\out{x}{\drp{z}}$ and both $x$
and $z$ are substitution-reachable, $z$ because
$\quo{(\drp{z})} \nequiv z$ makes $\drp{z}$ a process position and not a
quote.

$P = x(w).P_1$: the image is $\inp{w}{x}{\enc{P_1}_\sigma}$. If $y = w$
there is nothing to prove; otherwise substitution descends through the
input constructor into the continuation and the hypothesis applies.

$P = P_1 \mid P_2$: the four components of the clause are $\nil$-free
processes built from $\sigma$ and the two sub-images, each of which
occurs as the continuation of an input --- a process position. Descend.

$P = (\nu z)P_1$: the image is $\Frsh(\sigma;z,m).\enc{P_1}_m$, which
unfolds to two nested inputs and two lifts whose objects mention only
$\sigma$; $y \neq z$ by $\alpha$-conversion, and the hypothesis applies to
the continuation.

$P = \,!\,x(w).P_1$: the image is
$\Frsh(\sigma;a,b).\Frsh(b;s,c).(\Dup(a) \mid \out{a}{B} \mid \out{s}{\drp{c}})$
with $B$ as in Definition~\ref{def:translation}. The occurrences of $y$
are in $x$, if $y = x$, and in $B$. The former is the subject of the
input heading $B$; the latter is inside $B$, which is the object of the
lift $\out{a}{B}$, and by Lemma~\ref{lem:transparency}
$(\out{a}{B})\{n/y\} = \out{a}{(B\{n/y\})}$. The hypothesis applies to
$\enc{P_1}_d$ inside $B$.

The side condition on $n$ is what keeps $\rho[y\mapsto n]$ an address
environment.
\end{proof}

Proposition~\ref{prop:subst} is where admissibility
(Definition~\ref{def:discipline}) is used a second time: it is exactly the
statement that the translation writes no parameter under a quote, and it
fails for the $2005$ clauses by inspection.

\subsection{B.2 Administrative and principal reductions}

\begin{definition}\label{def:phases}
A reduction of a term reachable from $\enc{P}^\rho_{\sigma}$ is
\emph{administrative} if the subject of its redex is an address, and
\emph{principal} if the subject is in $\varphi(\fn(P))$ or in the image of
an address environment --- that is, if it is a source name. Write
$\to_a$ and $\to_p$.
\end{definition}

\begin{lemma}\label{lem:exhaustive}
Every reduction of a reachable term is administrative or principal, and
not both.
\end{lemma}

\begin{proof}
Exhaustiveness is Theorem~\ref{thm:locality}(i). Exclusivity: an address
is $\theta(n_0)$ for $\theta \in \{\Lft,\Rgt\}^*$ and $n_0$ is chosen
outside $\varphi(\fn(P))$; by Lemma~\ref{lem:templates} the map
$\theta \mapsto \theta(n_0)$ is injective and its image contains no
$\varphi(x)$. The names in the image of an address environment are
addresses, but they are addresses \emph{allocated to a $\nu$}, and by
Theorem~\ref{thm:locality}(ii) an address allocated to a $\nu$ site is the
subject of no clause's bookkeeping; it is used only as a source name.
\end{proof}

\begin{lemma}[Commutation]\label{lem:commute}
If $T \to_a T_1$ and $T \to_p T_2$ then there is $T_3$ with
$T_1 \to_p T_3$ and $T_2 \to_a T_3$.
\end{lemma}

\begin{proof}
By Lemma~\ref{lem:exhaustive} the two redexes have subjects that are not
$\nequiv$-equal, so they share neither their input nor their lift, and are
therefore disjoint occurrences in $T$ up to $\equiv$. Reduction is closed
under $\mid$ and $\equiv$, and neither step alters the other's redex:
the principal step substitutes a name for a variable bound by a source
input, and by Lemma~\ref{lem:usage} that variable occurs in no address
position; the administrative step substitutes an address for a variable
bound by a translation input, which by construction does not occur in the
principal redex.
\end{proof}

\begin{corollary}[Phase separation]\label{cor:phases}
Every reduction sequence $T \to^{k} T'$ may be reordered, without
changing $T'$, so that each principal step is preceded by all the
administrative steps available before it.
\end{corollary}

\begin{lemma}[The seed is always returned]\label{lem:seed}
In an unfolding of a replication node, the message $\out{s}{\drp{e}}$
returning the successor seed is emitted after two administrative
reductions on $c'$ and depends on no other communication. Consequently a
replication node with several unfoldings in progress serialises them on
$s$ without deadlock, and every maximal administrative sequence from a
reachable term is finite.
\end{lemma}

\begin{proof}
The continuation of the seed receipt is
$\Frsh(c';d,e).(\out{s}{\drp{e}} \mid \enc{P}_d)$, and $\Frsh$ supplies its
own two messages on $c'$; the two receipts are therefore enabled
immediately and the emission of $\out{s}{\drp{e}}$ is guarded by nothing
else. Finiteness is then Lemma~\ref{lem:admin}, whose count is per
unfolding and whose unfoldings are guarded by principal steps.
\end{proof}

\subsection{B.3 Realisation}

\begin{definition}[Realisation]\label{def:real}
Let $\real$ be the least relation between rho processes and
configurations such that
\begin{enumerate}[label=(\roman*),leftmargin=2.2em,nosep]
\item $\enc{P}^\rho_\sigma \real (P,\rho,\sigma)$;
\item if $T \real (P,\rho,\sigma)$ and $T \to_a T'$ then
$T' \real (P,\rho',\sigma)$, where $\rho' = \rho[z \mapsto \alpha]$ if the
step is the second receipt of a $\Frsh$ occupying the position of a
$\nu$-bound $z$, and $\rho' = \rho$ otherwise;
\item if $T \real (P,\rho,\sigma)$ and $T \equiv T'$ then
$T' \real (P,\rho,\sigma)$.
\end{enumerate}
\end{definition}

Clause (ii) is the point of the definition: administrative work does not
change which source term is being realised, but it may discharge a $\nu$,
and that is recorded in the environment rather than in the source term.

\begin{lemma}[Realisation is compositional]\label{lem:real-comp}
If $T \real (P,\rho,\sigma)$ and $C$ is a $\pi$ context with
$\enc{C}$ its image, then $\enc{C}[T] \real (C[P],\rho,\sigma_C)$ for the
address $\sigma_C$ at which $\enc{C}$ invokes its hole.
\end{lemma}

\begin{proof}
By Lemma~\ref{lem:compositional} the base case (i) is
$\enc{C}[\enc{P}^\rho_\sigma] = \enc{C[P]}^\rho_{\sigma_C}$. Clauses (ii)
and (iii) are preserved because reduction and $\equiv$ are closed under
parallel composition and under the prefixes the translation builds, and
because the addresses $\enc{C}$ allocates are disjoint from those below
$\sigma$ by Theorem~\ref{thm:locality}(ii).
\end{proof}

The structural content of the appendix is the following lemma, which says
that once the bookkeeping is finished the term exposes exactly the source
term's own capabilities and nothing else.

\begin{lemma}[Exposure]\label{lem:exposure}
Let $T \real (P,\rho,\sigma)$ with $T$ administratively normal. Then
the unguarded prefixes of $T$ whose subject is a source name are in
bijection with the unguarded prefixes of $P$, the bijection sending a
$\pi$ prefix with subject $x$ to a rho prefix with subject $\hat\rho(x)$,
and:
\begin{enumerate}[label=(\roman*),leftmargin=2.2em,nosep]
\item an unguarded $\bar x\langle z\rangle$ of $P$ corresponds to
$\out{\hat\rho(x)}{\drp{\hat\rho(z)}}$;
\item an unguarded $x(y).P_1$ of $P$ corresponds to
$\inp{y}{\hat\rho(x)}{U}$ with $U \real (P_1,\rho,\sigma_1)$ for the
address $\sigma_1$ assigned to that occurrence;
\item an unguarded $!\,x(y).P_1$ of $P$ corresponds to a pair
$\out{a}{B} \mid B$ with $B = \inp{y}{\hat\rho(x)}{\cdots}$ as in
Definition~\ref{def:translation}, and $a$ the box channel allocated to
that occurrence.
\end{enumerate}
Moreover every unguarded prefix of $T$ whose subject is an address has no
matching co-prefix, $T$ being administratively normal.
\end{lemma}

\begin{proof}
Induction on $P$, with the base cases (i)--(iii) read off
Definition~\ref{def:translation}.

$\nil$, output and input allocate nothing, pass their address to the
continuation, and their images are already administratively normal;
(i) and (ii) hold by inspection.

$P_1 \mid P_2$: the clause offers two messages on $\sigma$ and two
receivers on $\sigma$, so exactly two administrative reductions are
available and after them $T \equiv \enc{P_1}^\rho_{m_1} \mid
\enc{P_2}^\rho_{m_2}$ with $\{m_1,m_2\} = \{\Lft(\sigma),\Rgt(\sigma)\}$
in one of the two assignments. Both assignments satisfy the statement,
and they are related by $\beta$ exchanging the two subtrees
(Corollary~\ref{cor:race}). Apply the hypothesis to each factor; the
address subtrees are disjoint by Lemma~\ref{lem:templates}, so the
bijections combine.

$(\nu z)P_1$: two administrative reductions discharge the $\Frsh$. They
produce $\enc{P_1}^{\rho'}_{\beta}$, where
$\rho' = \rho[z\mapsto\alpha]$ and
$\{\alpha,\beta\} = \{\Lft(\sigma),\Rgt(\sigma)\}$, and clause (ii) of
Definition~\ref{def:real} has recorded $z$. Apply the hypothesis. Note
that $\alpha$ is thereafter a source name in the sense of
Definition~\ref{def:phases}, and is by
Theorem~\ref{thm:locality}(ii) the subject of no bookkeeping.

$!\,x(y).P_1$: four administrative reductions discharge the two nested
$\Frsh$, giving $\Dup(a) \mid \out{a}{B} \mid \out{s}{\drp{c_0}}$; a
fifth, the copy $\Dup(a) \mid \out{a}{B} \to B \mid \out{a}{B}$, is the
last available, since $B$ is guarded by an input on the source name $x$
and no duplicator remains on $a$. The message $\out{s}{\drp{c_0}}$ has no
receiver until $B$ fires. Hence the administrative normal form is
$\out{a}{B} \mid B \mid \out{s}{\drp{c_0}}$, which is (iii), and the
remaining unguarded prefix with an address subject, namely
$\out{s}{\drp{c_0}}$, has no co-prefix as claimed.

For the induction to close we must also treat a $T$ obtained by clause
(ii) after a principal step, that is a term realising a $P$ with some
replication node partly unfolded. Such a node contributes
$\out{a}{B} \mid B \mid \out{s}{\drp{c_k}} \mid \enc{P_1\{z/y\}}^\rho_{d_k}$
once its administrative work is finished, by Lemma~\ref{lem:seed}, and the
first three conjuncts are again (iii) while the fourth is covered by the
hypothesis. This is the case in which a box has been copied and the seed
exchange not yet performed: it is administratively reducible, hence
excluded from the statement, and Lemma~\ref{lem:seed} says the exchange
always completes.
\end{proof}

\subsection{B.4 Completeness}

\begin{proof}[Proof of Proposition~\ref{prop:complete}]
Let $P \to P'$. By definition of $\pi$ reduction, $P \equiv
C[\bar x\langle z\rangle \mid x(y).P_1]$ for a $\pi$ context $C$ and $P'
\equiv C[P_1\{z/y\}]$, or the same with $!\,x(y).P_1$ in place of
$x(y).P_1$ and $P' \equiv C[\,!\,x(y).P_1 \mid P_1\{z/y\}\,]$.

Take $T$ the administrative normal form of $\enc{P}^\rho_\sigma$, which
exists by Lemma~\ref{lem:seed} and is unique up to $\beta$ by
Corollary~\ref{cor:race} and Lemma~\ref{lem:addrind}. By
Lemma~\ref{lem:exposure} the two prefixes of the source redex appear in
$T$ as $\out{\hat\rho(x)}{\drp{\hat\rho(z)}}$ and either
$\inp{y}{\hat\rho(x)}{U}$ with $U \real (P_1,\rho,\sigma_1)$, or the box
$\out{a}{B} \mid B$. Since $\quo{(\drp{\hat\rho(z)})} \nequiv
\hat\rho(z)$, one principal reduction gives, in the first case,
$U\{\hat\rho(z)/y\}$, which by Proposition~\ref{prop:subst} is
$\enc{P_1\{z/y\}}^\rho_{\sigma_1}$ up to the administrative work already
performed inside $U$, hence realises $(P_1\{z/y\},\rho,\sigma_1)$; and in
the second case
\[
  \out{a}{B} \ \mid\ \Dup(a) \ \mid\
  \inp{c'}{s}{\Frsh(c';d,e).(\out{s}{\drp{e}} \mid \enc{P_1\{z/y\}}_d)}
\]
which by Lemma~\ref{lem:seed} administratively reduces to
$\out{a}{B} \mid B \mid \out{s}{\drp{c_{k+1}}} \mid
\enc{P_1\{z/y\}}^\rho_{d_k}$, realising
$(\,!\,x(y).P_1 \mid P_1\{z/y\},\,\rho,\,\sigma\,)$.

In both cases the surrounding context contributes
$\enc{C}$ and Lemma~\ref{lem:real-comp} assembles the result, so
$\enc{P}^\rho_\sigma \to^{+} T'$ with $T' \real (P',\rho,\sigma)$, whence
$T' \bisim_{\enc{\Ctx_\pi}} \enc{P'}^\rho_{\sigma'}$ for the address
$\sigma'$ of the corresponding occurrence, by
Lemma~\ref{lem:addrind}. Reduction closed under $\equiv$ handles the
structural congruence in the first line; the case $P \mid \nil \equiv P$
uses Lemmas~\ref{lem:admin} and \ref{lem:addrind} to discard the
administrative residue of the $\nil$ branch, which by
Theorem~\ref{thm:locality}(i) is unobservable on $\varphi(\fn(P))$.
\end{proof}

\subsection{B.5 Soundness}

\begin{proof}[Proof of Proposition~\ref{prop:sound}]
Let $\enc{P}^\rho_\sigma \to^{k} T$. By Corollary~\ref{cor:phases} we may
assume the sequence is $A_0\,p_1\,A_1\,p_2 \cdots p_j\,A_j$ with each
$A_i$ a maximal administrative block and each $p_i$ principal; the blocks
are finite by Lemma~\ref{lem:seed}. Induct on $j$.

$j = 0$: $T \real (P,\rho,\sigma)$ by clause (ii) of
Definition~\ref{def:real}, and $P \to^{0} P$; the required
$T \bisim_{\enc{\Ctx_\pi}} \enc{P}^{\rho}_{\sigma'}$ follows because
administrative steps are unobservable on $\varphi(\fn(P))$
(Lemma~\ref{lem:admin}) and confluent up to $\beta$
(Corollary~\ref{cor:race}), so the bisimulation
$\{(T_1,T_2) : T_1,T_2 \real (P,\rho,\sigma)\}$ relates any two states of
one configuration; it is a $\enc{\Ctx_\pi}$-bisimulation because an
encoded context can neither observe nor participate in a communication
whose subject is an address, by Theorem~\ref{thm:locality}(ii) and
Lemma~\ref{lem:usage}.

$j > 0$: let $T_0$ be the term at the end of $A_0$, so $T_0$ is
administratively normal and $T_0 \real (P,\rho_0,\sigma)$. The step $p_1$
is principal, so its subject is a source name, and by
Lemma~\ref{lem:exposure} its redex consists of two prefixes that are the
images of unguarded prefixes of $P$ with a common subject. Hence
$P \to P_1$ by the corresponding source communication, and the argument of
B.4 read backwards identifies the resulting term as realising
$(P_1,\rho_0,\sigma)$. Apply the hypothesis to the remaining
$j-1$ principal steps.

The classification is exhaustive by Lemma~\ref{lem:exhaustive}; the only
non-confluence among administrative steps is the $\Frsh$ race, absorbed by
Corollary~\ref{cor:race}; and the case in which a box has been copied but
the seed exchange has not occurred is a term in the middle of a block
$A_i$, which the reordering of Corollary~\ref{cor:phases} places before
the next principal step and Lemma~\ref{lem:seed} guarantees completes.
\end{proof}

\begin{remark}
Soundness in the form above is the ``operational correspondence, converse
direction'' of \cite{Gorla10} and is where an encoding usually pays for
administrative steps. The price here is bounded: by
Lemma~\ref{lem:admin} the number of administrative steps between two
principal ones is a function of the source node performing them, and by
Theorem~\ref{thm:onename} each of them buys at most one name.
\end{remark}

\subsection{B.6 Full abstraction}

We use the standard reformulation of $\bisim_\Obs$ in which a single
labelled step is matched by a labelled path; that the two definitions
agree is proved as for weak bisimilarity, by transitivity of $\to^*$.

\begin{lemma}[Weak transfer]\label{lem:weaktransfer}
A symmetric $\mathcal{R}$ is a $\Obs$-bisimulation if and only if
$P \mathrel{\mathcal{R}} Q$ and $P \xLongrightarrow{\,C\,} P'$ with
$C \in \Obs$ imply $Q \xLongrightarrow{\,C\,} Q'$ with
$P' \mathrel{\mathcal{R}} Q'$.
\end{lemma}

\begin{proof}[Proof of Theorem~\ref{thm:fullabs}]
\emph{Forward.} Let
\[
  \mathcal{S} \ =\ \{(T,U) \ :\ \exists P,Q,\rho,\sigma.\
  P \bisim_\pi Q,\ T \real (P,\rho,\sigma),\ U \real (Q,\rho,\sigma)\}.
\]
$\mathcal{S}$ is symmetric and contains $(\enc{P},\enc{Q})$ whenever
$P \bisim_\pi Q$, by clause (i) of Definition~\ref{def:real}. We check it
is a $\enc{\Ctx_\pi}$-bisimulation, in the form of
Lemma~\ref{lem:weaktransfer}.

Let $T \mathrel{\mathcal{S}} U$ and $\enc{C}[T] \to^* T'$ for a $\pi$
context $C$. By Lemma~\ref{lem:real-comp}, $\enc{C}[T] \real
(C[P],\rho,\sigma_C)$. By Proposition~\ref{prop:sound} there is $P'$ with
$C[P] \to^* P'$ and $T' \real (P',\rho',\sigma')$. Since
$P \bisim_\pi Q$ and $C$ is a source context, $C[Q] \Longrightarrow^{\,C\,}$
matches: there is $Q'$ with $C[Q] \to^* Q'$ and $P' \bisim_\pi Q'$. By
Proposition~\ref{prop:complete} applied along that source path,
$\enc{C}[U] \to^* U'$ with $U' \real (Q',\rho',\sigma'')$. Then
$(T',U') \in \mathcal{S}$, using Lemma~\ref{lem:addrind} to absorb the
difference between $\sigma'$ and $\sigma''$.

\emph{Backward.} Let
$\mathcal{T} = \{(P,Q) : \enc{P} \bisim_{\enc{\Ctx_\pi}} \enc{Q}\}$ and
suppose $C[P] \to^* P'$ for a $\pi$ context $C$. By
Proposition~\ref{prop:complete}, $\enc{C}[\enc{P}] \to^* T'$ with
$T' \real (P',\rho,\sigma)$. By hypothesis and
Lemma~\ref{lem:weaktransfer} there is $U'$ with
$\enc{C}[\enc{Q}] \to^* U'$ and $T' \bisim_{\enc{\Ctx_\pi}} U'$. By
Proposition~\ref{prop:sound} there is $Q'$ with $C[Q] \to^* Q'$ and
$U' \real (Q',\rho',\sigma')$. It remains to see
$(P',Q') \in \mathcal{T}$, that is
$\enc{P'} \bisim_{\enc{\Ctx_\pi}} \enc{Q'}$; this holds because
$T' \bisim_{\enc{\Ctx_\pi}} \enc{P'}$ and
$U' \bisim_{\enc{\Ctx_\pi}} \enc{Q'}$, each by the base case of the
soundness proof, and $\bisim_{\enc{\Ctx_\pi}}$ is transitive. Hence
$\mathcal{T}$ is a $\pi$ bisimulation and $P \bisim_\pi Q$.
\end{proof}

\begin{remark}
The backward direction is where the restriction of observers is
indispensable, and it is instructive to see exactly where. Nothing above
would fail if $\Obs$ were larger, \emph{except} the base case of the
soundness proof, which asserts that two states of one configuration are
indistinguishable. An unrestricted rho context distinguishes them ---
Proposition~\ref{prop:necessity} exhibits a pair --- and with that step
gone the chain in the backward direction breaks at its first link. The
restriction is not used to make the theorem easier; it is used once, and
in the one place where the encoding's bookkeeping becomes visible.
\end{remark}

\subsection{B.7 Lybech's criteria}

\begin{proof}[Proof of Proposition~\ref{prop:criteria}]
Take $\simeq$ to be $\bisim_{\enc{\Ctx_\pi}}$ and the name derivation
function to be $\sigma \mapsto \{\Lft(\sigma),\Rgt(\sigma)\}$.
\emph{Compositionality} is Lemma~\ref{lem:compositional}, which holds on
the nose rather than up to a context, because $\enc{P}$ is an address
abstraction and plugging is instantiation.
\emph{Name invariance} and \emph{substitution invariance} are
Proposition~\ref{prop:subst} together with Lemma~\ref{lem:usage}: a
permutation of source names commutes with the translation, since source
names occur only in subject and lift-object positions.
\emph{Operational correspondence} is
Propositions~\ref{prop:complete} and \ref{prop:sound}.
\emph{Divergence reflection} is the corollary of
Lemma~\ref{lem:admin}: an infinite target reduction sequence contains
infinitely many principal steps, since administrative blocks are finite,
and each principal step is a source communication.
\emph{Observational correspondence} follows from
Theorem~\ref{thm:locality}(i): the barbs of $\enc{P}$ on
$\varphi(\fn(P))$ are the images of the barbs of $P$, no administrative
step being observable there.
\emph{Parameter independence} is Lemma~\ref{lem:addrind}, and is the
criterion the $2005$ encoding fails.
\end{proof}

\begin{remark}
Only the last of these is proved differently here than it would be for a
server-based encoding. Lybech obtains parameter independence as a
property of his translation; here it is a corollary of address
relabelling, which is a symmetry of the calculus restricted to the
encoded contexts, and it is proved once for the whole development rather
than re-established per clause. That is the practical dividend of
Definition~\ref{def:discipline}.
\end{remark}

% =====================================================================
\section{The encoding in rholang}\label{app:rholang}
% =====================================================================

\begin{lstlisting}
// Addresses:  L(s) = @{ s!(Nil) }        R(s) = @{ for (y <- s) { Nil } }

// Fresh2(s; u, v).P   --- monadic form, as in the calculus:
//   for (u <- s) { for (v <- s) { P } }
//   | s!( s!(Nil) ) | s!( for (y <- s) { Nil } )
//
// Fresh2(s; u, v).P   --- join form, as in rholang (deterministic):
//   for (u <- s & v <- s) { P }
//   | s!( s!(Nil) ) | s!( for (y <- s) { Nil } )

// [[ P1 | P2 ]]_s
//   for (m <- s) { [[P1]]_m } | for (m <- s) { [[P2]]_m }
//   | s!( s!(Nil) ) | s!( for (y <- s) { Nil } )

// [[ (nu z) P ]]_s
//   Fresh2(s; z, m). [[P]]_m

// [[ ! x(y).P ]]_s
//   Fresh2(s; a, b). Fresh2(b; srv, c).
//     ( D(a) | a!( B ) | srv!( *c ) )
//   B = for (y <- x) {
//         D(a) |
//         for (cp <- srv) {
//           Fresh2(cp; d, e).( srv!(*e) | [[P]]_d )
//         }
//       }
//   D(a) = for (q <- a) { *q | a!(*q) }
\end{lstlisting}

\noindent
The join form of $\Frsh$ consumes both messages in one rendezvous and so
has no race; it is what an implementation should use, and it is the
concrete payoff of \S\ref{sec:notation}(2). The theory above assumes only
the monadic form, and pays for that assumption in
Corollary~\ref{cor:race}.

\end{document}
